% Luc Destombe --- licence CC BY-NC-SA 4.0
% https://creativecommons.org/licenses/by-nc-sa/4.0/deed.fr
% Fichier autonome : compiler avec pdflatex, deux fois (signets, nombre de pages).
\documentclass[a4paper,10pt]{article}
\makeatletter
\newif\iflycee@niveaux
\lycee@niveauxtrue
\newif\iflycee@palatino
\lycee@palatinotrue

\RequirePackage[utf8]{inputenc}
\RequirePackage[T1]{fontenc}
\RequirePackage[french]{babel}
\RequirePackage{etoolbox}

\newcommand{\ShorthandsOff}{\@ifundefined{shorthandoff}{}{\shorthandoff{;:!?}}}
\newcommand{\ShorthandsOn}{\@ifundefined{shorthandon}{}{\shorthandon{;:!?}}}
\ShorthandsOff
\AtBeginDocument{\ShorthandsOn}
\AtBeginEnvironment{tikzpicture}{\ShorthandsOff}

\RequirePackage{geometry}
\RequirePackage{fancyhdr}
\RequirePackage{lastpage}
\RequirePackage{multicol}
\RequirePackage{array}
\RequirePackage{enumitem}
\RequirePackage{xcolor}
\RequirePackage{needspace}

\RequirePackage{amsmath,amssymb}
\RequirePackage{mathpazo}
\DeclareMathSizes{10.95}{10.95}{8}{5.5}
\begingroup\catcode`\_=\active \catcode`\^=\active
\gdef_#1{\sb{\scriptscriptstyle #1}}
\gdef^#1{\sp{\scriptscriptstyle\lycee@moinsepais #1}}
\endgroup
\newcommand\lycee@moins{\mathbin{%
  \setbox\z@\hbox{$\m@th\scriptscriptstyle\lycee@moinsorig$}%
  \dimen@=.55\wd\z@ \advance\dimen@-.5pt
  \hbox to.75\wd\z@{\hss$\m@th\scriptscriptstyle\vcenter{\hbox{%
    \tikz\draw[line width=.5pt,line cap=round](0,0)--(\the\dimen@,0);}}$\hss}}}
\begingroup\lccode`\~=`\- \lowercase{\endgroup\let~}\lycee@moins
\newcommand\lycee@moinsepais{\mathcode`\-="8000 }
\AtBeginDocument{\mathchardef\lycee@moinsorig=\mathcode`\-\relax}
\AtBeginDocument{\catcode`\_=12 \mathcode`\_="8000
  \catcode`\^=12 \mathcode`\^="8000 }
\newcommand\lycee@exposants{%
  \fontdimen13\textfont2=.48\fontdimen6\textfont2
  \fontdimen14\textfont2=.48\fontdimen6\textfont2
  \fontdimen15\textfont2=.48\fontdimen6\textfont2
  \fontdimen13\scriptfont2=.48\fontdimen6\scriptfont2
  \fontdimen14\scriptfont2=.48\fontdimen6\scriptfont2
  \fontdimen15\scriptfont2=.48\fontdimen6\scriptfont2}
\every@math@size\expandafter{\the\every@math@size\lycee@exposants}
\newlength{\retraitcalculs}
\setlength{\retraitcalculs}{2em}

\RequirePackage{tikz}
\usetikzlibrary{arrows.meta,calc,babel,shapes.misc}
\RequirePackage[most]{tcolorbox}
\tcbuselibrary{skins,breakable}

\definecolor{coulDef}{HTML}{1F4E79}
\definecolor{coulProp}{HTML}{7030A0}
\definecolor{coulRem}{HTML}{C55A11}
\definecolor{coulEx}{HTML}{2E7D32}
\definecolor{coulExo}{HTML}{B03A2E}
\definecolor{coulPart}{HTML}{1F4E79}
\definecolor{coulMeth}{HTML}{1B6E4F}
\definecolor{coulCourbe}{HTML}{0B5ED7}
\definecolor{coulCourbeB}{HTML}{C00000}
\definecolor{coulCourbeC}{HTML}{2E7D32}
\definecolor{coulGrille}{HTML}{8C8C8C}
\definecolor{coulRep}{HTML}{1B6E4F}
\definecolor{coulBar}{HTML}{2E7D32}

\newcommand{\R}{\mathbb{R}}
\newcommand{\Rs}{\mathbb{R}^{*}}
\renewcommand{\le}{\leqslant}
\renewcommand{\ge}{\geqslant}
\newcommand{\vect}[1]{\overrightarrow{#1}}
\newcommand{\norme}[1]{\left\lVert #1 \right\rVert}
\newcommand{\intervalle}[2]{\left[#1\,;#2\right]}
\RequirePackage{cancel}
\renewcommand{\CancelColor}{\color{coulExo}}
\newcommand{\fois}[1]{\times{\color{coulExo}#1}}
\newcommand{\barre}[1]{\cancel{#1}}

\tikzset{grille/.style={coulGrille,line width=0.4pt}}
\newcommand{\quadrillage}[4]{\draw[grille,step=1] (#1,#2) grid (#3,#4);}
\tikzset{
  croix/.style={cross out,draw,line width=0.6pt,inner sep=0pt,minimum size=4pt},
  croixdroite/.style={croix,rotate=45,line width=0.8pt},
}
\newcommand{\pt}[3]{\path (#1) node[croix]{} node[#2,font=\small,inner sep=2.5pt]{$#3$};}

\newcommand{\lycee@repere}[4]{%
  \draw[grille,step=1] (#1,#2) grid (#3,#4);
  \draw[-{Stealth[length=2mm]},thick] (#1,0)--({#3+0.4},0) node[below left=-1pt]{$x$};
  \draw[-{Stealth[length=2mm]},thick] (0,#2)--(0,{#4+0.4}) node[below left=-1pt]{$y$};
  \node[below left,font=\scriptsize,inner sep=1.5pt] at (0,0) {$0$};}
\newcommand{\lycee@gradx}[1]{\foreach \k/\l in {#1}{%
  \draw (\k,0.07)--(\k,-0.07) node[below,font=\scriptsize,inner sep=1.5pt]{$\l$};}}
\newcommand{\lycee@grady}[1]{\foreach \k/\l in {#1}{%
  \draw (0.07,\k)--(-0.07,\k) node[left,font=\scriptsize,inner sep=1.5pt]{$\l$};}}
\newcommand{\lycee@intff}[2]{\left[#1\,;#2\right]}
\newcommand{\lycee@intfo}[2]{\left[#1\,;#2\right[}
\newcommand{\lycee@intof}[2]{\left]#1\,;#2\right]}
\newcommand{\lycee@intoo}[2]{\left]#1\,;#2\right[}
\newcommand{\lycee@ens}[1]{\left\{#1\right\}}
\AtBeginDocument{%
  \providecommand{\repere}{\lycee@repere}%
  \providecommand{\gradx}{\lycee@gradx}%
  \providecommand{\grady}{\lycee@grady}%
  \providecommand{\Cf}{\mathcal{C}\sb{\scriptscriptstyle f}}%
  \providecommand{\Cg}{\mathcal{C}\sb{\scriptscriptstyle g}}%
  \providecommand{\intff}{\lycee@intff}%
  \providecommand{\intfo}{\lycee@intfo}%
  \providecommand{\intof}{\lycee@intof}%
  \providecommand{\intoo}{\lycee@intoo}%
  \providecommand{\ens}{\lycee@ens}}

\newlist{questions}{enumerate}{3}
\setlist[questions,1]{label=\arabic*\textdegree),leftmargin=*,itemsep=2pt,topsep=3pt}
\setlist[questions,2]{label=\alph*),leftmargin=*,itemsep=1pt,topsep=2pt}
\setlist[questions,3]{label=\roman*),leftmargin=*,itemsep=1pt,topsep=1pt}
\setlist[itemize]{label=\textbullet,leftmargin=*,itemsep=2pt,topsep=3pt}

\newcommand{\besoinplace}[1]{\par\penalty\@M\begingroup
  \dimen@=#1\relax\dimen@ii=\pagegoal\advance\dimen@ii-\pagetotal
  \ifdim\dimen@>\dimen@ii\ifdim\dimen@ii>\z@\vfil\fi\break\fi\endgroup}

\newcommand{\lycee@pied}{\thepage/\pageref{LastPage}}
\newcommand{\entete}[2]{%
  \pagestyle{fancy}\fancyhf{}%
  \fancyhead[L]{\footnotesize\color{coulPart}\textbf{#1}}%
  \fancyhead[R]{\footnotesize\color{coulPart}\textbf{#2}}%
  \fancyfoot[C]{\footnotesize\lycee@pied}%
  \renewcommand{\headrulewidth}{0.4pt}%
  \renewcommand{\headrule}{\hbox to\headwidth{%
    \color{coulPart!40}\leaders\hrule height \headrulewidth\hfill}}}

\RequirePackage{ccicons}
\newcommand{\lycee@licence}{{\scriptsize\color{gray}%
  \copyright\ Luc Destombe --- \ccbyncsa\ CC BY-NC-SA 4.0}}
\newif\iflycee@licence \lycee@licencetrue
\newcommand{\sanslicence}{\lycee@licencefalse}
\AtBeginDocument{\iflycee@licence\AddToHookNext{shipout/foreground}{%
  \put(0,0){\rlap{\kern\dimexpr1in+\hoffset+\oddsidemargin+\textwidth\relax
    \llap{\raisebox{-\dimexpr1in+\voffset+\topmargin+\headheight+\headsep
      +\textheight+\footskip\relax}{\lycee@licence}}}}}\fi}

\newcommand{\formulesserrees}{%
  \setlength{\abovedisplayskip}{3pt}\setlength{\belowdisplayskip}{3pt}%
  \setlength{\abovedisplayshortskip}{2pt}\setlength{\belowdisplayshortskip}{3pt}}

\setlength{\parindent}{0pt}

\RequirePackage[implicit=false,hidelinks,pdfencoding=auto,psdextra,bookmarksopen,
  bookmarksopenlevel=1,pdfcreator={},pdfproducer={}]{hyperref}
\RequirePackage{bookmark}
\hypersetup{pdfauthor={Luc Destombe},pdfkeywords={CC BY-NC-SA 4.0}}
\newcounter{lycee@signet}
\newcommand{\signet}[3][]{\stepcounter{lycee@signet}%
  \edef\lycee@dest{\if\relax\detokenize{#1}\relax
    signet.\arabic{lycee@signet}\else#1\fi}%
  \hypertarget{\lycee@dest}{}\bookmark[level=#2,dest=\lycee@dest]{#3}}
\pdfstringdefDisableCommands{%
  \def\R{R}\def\vect#1{#1}\def\Cf{Cf}\def\Cg{Cg}\def\fois#1{×#1}%
  \def\barre#1{#1}\def\intervalle#1#2{[#1;#2]}\def\intff#1#2{[#1;#2]}%
  \def\textsuperscript#1{#1}\def\dfrac#1#2{#1/#2}\def\frac#1#2{#1/#2}%
  \def\vec#1{vecteur #1}}
\RequirePackage[official]{eurosym}

\geometry{top=0.8cm,bottom=1.3cm,left=1cm,right=1cm,footskip=0.6cm}

\pagestyle{fancy}
\fancyhf{}
\renewcommand{\headrulewidth}{0pt}
\fancyfoot[C]{\footnotesize Page \thepage{} / \pageref{LastPage}}

\newcommand{\titrefiche}[2]{%
  \begin{center}
    \Large\textbf{#1}\\[2pt]
    \large\textbf{#2}\\[2pt]
  \end{center}
  \vspace{0.10cm}}

\setlength{\columnseprule}{0.4pt}
\setlength{\columnsep}{15pt}
\renewcommand{\columnseprulecolor}{\color{black!45}}
\setlength{\multicolsep}{2pt}

\newcounter{partiefiche}
\renewcommand{\thepartiefiche}{\Roman{partiefiche}}
\newif\ifapresbandeau
\newcommand{\lycee@bandeau}[1]{%
  \par\penalty-600
  \vspace{0.08cm}%
  \noindent\colorbox{coulPart}{%
    \parbox{\dimexpr\linewidth-2\fboxsep\relax}{%
      \color{white}\bfseries\raggedright\strut#1\strut}}%
  \nopagebreak\par\nopagebreak
  \vspace{0.06cm}\nopagebreak
  \apresbandeautrue}
\newcommand{\partiefiche}[1]{%
  \refstepcounter{partiefiche}%
  \lycee@bandeau{\signet{1}{\thepartiefiche{} --- #1}\thepartiefiche{} --- #1}}
\newcommand{\partiefichesansnum}[1]{\lycee@bandeau{\signet{1}{#1}#1}}

\newcounter{exo}
\newlength{\espaceexo}\setlength{\espaceexo}{7pt}
\newlength{\espacetitre}\setlength{\espacetitre}{2pt}
\newcommand{\exo}[1][\relax]{%
  \par
  \ifapresbandeau\apresbandeaufalse\else\penalty-150\addvspace{\espaceexo}\fi
  \refstepcounter{exo}%
  \noindent\signet[exercice-\arabic{exo}]{2}{Exercice \arabic{exo}}%
  \textbf{\color{coulExo}\underline{Exercice \theexo}}%
  \ifx\relax#1\relax\else\textbf{ : #1}\fi
  \par\nopagebreak\vspace{\espacetitre}\nopagebreak
  \ignorespaces}

\setlist[enumerate]{nosep,leftmargin=*,itemsep=2pt,topsep=2pt}
\setlist[itemize]{nosep,leftmargin=*,topsep=2pt}
\newcommand{\choix}[3]{\par\hspace*{1em}%
  \makebox[0.3\linewidth][l]{a) #1}\makebox[0.3\linewidth][l]{b) #2}c) #3}
\newcommand{\deux}[2]{\par\hspace*{0.6em}%
  \makebox[0.48\linewidth][l]{#1}#2}
\newcommand{\trois}[3]{\par\hspace*{0.6em}%
  \makebox[0.32\linewidth][l]{#1}\makebox[0.32\linewidth][l]{#2}#3}

  \renewcommand{\exo}[1][\relax]{%
    \ifapresbandeau\apresbandeaufalse\else\par\penalty-150\fi
    \refstepcounter{exo}%
    \vspace{0.05cm}%
    \noindent\signet[exercice-\arabic{exo}]{2}{Exercice \arabic{exo}}%
    \textbf{\color{coulExo}\underline{Exercice \theexo}}%
    \ifx\relax#1\relax\else\textbf{ : #1}\fi%
    \par\nopagebreak
    \ignorespaces}
  \newcommand{\rappel}[1]{%
    \par\vspace{0.06cm}\nopagebreak
    \noindent\fcolorbox{black!35}{black!5}{%
      \parbox{\dimexpr\linewidth-2\fboxsep-2\fboxrule\relax}{%
        \small\textbf{Rappel.} #1}}%
    \par\vspace{0.06cm}\nopagebreak}
  \setlist[enumerate]{nosep,leftmargin=*}
  \setlist[itemize]{nosep,leftmargin=*}
  \setlength{\multicolsep}{3pt plus 1pt minus 1pt}
  \setlength{\premulticols}{10pt}
  \setlength{\postmulticols}{10pt}
  \newlist{expr}{enumerate}{1}
  \setlist[expr]{label={\Alph*\ $=$},nosep,leftmargin=*,itemsep=0pt}

\renewenvironment{center}{\par\addvspace{3pt}\centering}{\par\addvspace{3pt}}
\formulesserrees
\AtBeginDocument{\formulesserrees}
\clubpenalty=10000
\widowpenalty=10000
\displaywidowpenalty=10000
\brokenpenalty=10000
\setlength{\parskip}{0pt}

\RequirePackage{ragged2e}
\setlength{\RaggedRightRightskip}{0pt plus 1fil}
\AtBeginDocument{\RaggedRight\binoppenalty=10000 \relpenalty=10000 }
\g@addto@macro\@parboxrestore{\RaggedRight}
\makeatother
\geometry{top=0.8cm,bottom=1.3cm,left=1cm,right=1cm,footskip=0.7cm,headheight=12pt}

\begin{document}
\titrefiche{1\textsuperscript{re} spécialité --- Calcul littéral}{Fiche d'exercices du chapitre}

\begin{multicols}{2}

\partiefiche{Réduire, développer, identités remarquables}

\exo[Réduire et distribuer]
Réduire, ou développer puis réduire :
\begin{multicols}{2}
\begin{expr}
  \item $3x+8+6x$
  \item $4x^{2}-9x^{2}+2x^{2}$
  \item $6(3x-2)$
  \item $2x(5x-4)$
  \item $9(x-1)-4x$
  \item $-5x+3(2x+7)$
  \item $-(4-3x)$
  \item $-2x(1-4x)$
\end{expr}
\end{multicols}

\exo[Double distributivité et identités remarquables]
Développer et réduire :
\begin{multicols}{2}
\begin{expr}
  \item $(2x+3)(4x+1)$
  \item $(3x-4)(-2x+5)$
  \item $(5x+2)^{2}$
  \item $(4-3x)^{2}$
  \item $(3x+5)(3x-5)$
  \item $\left(x-\dfrac{1}{3}\right)^{2}$
  \item $3(x+2)(1-x)$
\end{expr}
\end{multicols}

\exo[Développer, réduire, ordonner]
Développer, réduire, puis ordonner selon les puissances décroissantes de $x$ :
\begin{expr}
  \item $(4x+1)(2x+3)-(5x-2)(x+4)$
  \item $(2x-5)^{2}-(3x+1)(x-2)$
  \item $\left(x^{2}-x+2\right)(3x+1)-(x+1)(x-1)^{2}$
  \item $(x+2)\left(x^{2}-2x+4\right)$
  \item $3x(2-x)-(x-4)(2x+3)$
\end{expr}

\partiefiche{Substituer une valeur}

\exo[Calculer une image]
\begin{enumerate}[label=\arabic*),topsep=0pt]
  \item Soit $P(x)=3x^{2}-4x-2$. Calculer :
    \begin{enumerate}[label=\alph*)]
      \item $P(1)$ \quad et \quad $P(-1)$
      \item $P\!\left(\dfrac{1}{2}\right)$ \quad et \quad $P\!\left(-\dfrac{2}{3}\right)$
    \end{enumerate}
  \item Soient $A=(x+2)(x-5)$ et $B=x^{2}+3x-10$.
    \begin{enumerate}[label=\alph*)]
      \item Calculer $A$ et $B$ pour $x=0$.
      \item Calculer $A$ et $B$ pour $x=1$.
      \item Peut-on affirmer que $A=B$ pour tout réel $x$ ?
    \end{enumerate}
\end{enumerate}

\exo[Substituer une racine carrée]
Pour chaque polynôme $P$, calculer la valeur exacte de $P(a)$, sous la forme la plus
simple :
\begin{enumerate}[label=\alph*)]
  \item $P(x)=x^{2}-5$ \quad et \quad $a=\sqrt{5}$
  \item $P(x)=2x^{2}+3x-1$ \quad et \quad $a=\sqrt{2}$
  \item $P(x)=(3x+1)(x-2)$ \quad et \quad $a=\sqrt{3}$
  \item $P(x)=(2x-1)^{2}+4x$ \quad et \quad $a=\sqrt{5}$
  \item $P(x)=x^{2}-2x+3$ \quad et \quad $a=1+\sqrt{2}$
\end{enumerate}

\exo[Racines carrées dans un quotient]
\begin{enumerate}[label=\arabic*),topsep=0pt]
  \item Soit $f(x)=\dfrac{x+1}{x-3}$.
    \begin{enumerate}[label=\alph*)]
      \item Donner la valeur interdite de $f$.
      \item Calculer $f(5)$ et $f\!\left(\dfrac{1}{2}\right)$.
      \item Calculer $f\!\left(\sqrt{5}\right)$, sans racine carrée au dénominateur.
    \end{enumerate}
  \item Soit $g(x)=\dfrac{x-2}{x+1}$. Calculer $g\!\left(\sqrt{3}\right)$, sans racine
    carrée au dénominateur.
  \item Soit $h(x)=\dfrac{x+3}{x-2}$. Calculer $h\!\left(2+\sqrt{3}\right)$.
    \textit{Le dénominateur est ici une racine seule : on multiplie le numérateur et le
    dénominateur par cette racine.}
\end{enumerate}

\partiefiche{Factoriser une expression}

\exo[Facteur commun]
Factoriser :
\begin{multicols}{2}
\begin{expr}
  \item $5x+20$
  \item $7x^{2}+3x$
  \item $18x^{2}-12x$
  \item $x^{2}+x$
  \item $4x^{2}-x$
  \item $2x+x(x+5)$
\end{expr}
\end{multicols}

\exo[Identités remarquables et facteur commun]
Factoriser :
\begin{expr}
  \item $9x^{2}+30x+25$
  \item $49x^{2}-36$
  \item $(3x+1)(2x-5)+(3x+1)(x+7)$
  \item $(4x-1)(x+6)-(4x-1)(3x-2)$
  \item $(3x+2)^{2}-(x-4)^{2}$
  \item $36-(x+2)^{2}$
\end{expr}

\exo[Faire apparaître un facteur commun]
Factoriser. \textit{Pour A, B et C, factoriser d'abord une partie de l'expression.}
\begin{expr}
  \item $9x^{2}-16-(3x+4)(x-2)$
  \item $x^{2}-25-(x+5)(2x-3)$
  \item $(x+3)(4x-1)-\left(x^{2}-9\right)$
  \item $x^{2}-7$ \quad \textit{(écrire $7=\left(\sqrt{7}\right)^{2}$)}
  \item $x^{2}+2\sqrt{5}\,x+5$
\end{expr}

\partiefiche{Équations}

\exo[Équations du premier degré]
Résoudre dans $\R$ :
\begin{multicols}{2}
\begin{enumerate}[label=\alph*)]
  \item $4x-9=3x+2$
  \item $6x-5=x+20$
  \item $3x+4=3x-1$
  \item $3(x-2)=2(-x+7)$
  \item $x-(2-5x)=4x+9$
  \item $4(x-1)+3(x+5)=0$
\end{enumerate}
\end{multicols}

\exo[Produit nul et carrés]
Résoudre dans $\R$ :
\begin{multicols}{2}
\begin{enumerate}[label=\alph*)]
  \item $(x-6)(4+x)=0$
  \item $5x(3x-7)=0$
  \item $(2x+9)^{2}=0$
  \item $x^{2}=13$
  \item $3x^{2}-21=0$
  \item $(3x+1)^{2}=25$
\end{enumerate}
\end{multicols}

\exo[Mise en équation]
\textit{Pour chaque problème : choisir l'inconnue, écrire l'équation, la résoudre,
vérifier, puis conclure par une phrase.}
\begin{enumerate}[label=\arabic*)]
  \item \textbf{La boulangerie.} Un matin, une boulangerie a vendu $150$ viennoiseries :
    des croissants, des pains au chocolat et des chaussons aux pommes. Elle a vendu deux
    fois plus de pains au chocolat que de croissants, et $10$ chaussons de moins que de
    croissants. Combien de viennoiseries de chaque sorte a-t-elle vendues ?
  \item \textbf{Du carré au rectangle.} Un carré a pour côté $x$ (en cm). On allonge l'un
    de ses côtés de $6$ cm et on raccourcit l'autre de $4$ cm : le rectangle obtenu a la
    même aire que le carré. Trouver $x$, puis les dimensions du rectangle.
\end{enumerate}

\partiefiche{Fractions rationnelles}

\exo[Les quatre opérations]
Donner la valeur interdite, puis écrire chaque expression sous la forme d'un seul
quotient :
\begin{multicols}{2}
\begin{expr}
  \item $\dfrac{3}{4}+\dfrac{5}{x}$
  \item $\dfrac{3}{4}-\dfrac{5}{x}$
  \item $\dfrac{3}{4}\times\dfrac{5}{x}$
  \item $\dfrac{3}{4}\div\dfrac{5}{x}$
\end{expr}
\end{multicols}

\exo[Valeurs interdites]
Donner les valeurs interdites, puis simplifier quand c'est possible :
\begin{multicols}{2}
\begin{expr}
  \item $\dfrac{x+5}{x-4}$
  \item $\dfrac{6}{x(x-2)}$
  \item $\dfrac{x-4}{2x+5}$
  \item $\dfrac{x^{2}-16}{x-4}$
  \item $\dfrac{5}{x^{2}+4}$
  \item $\dfrac{2x+1}{7-x}$
\end{expr}
\end{multicols}

\exo[Réduire au même dénominateur]
Écrire sous la forme d'un seul quotient, en donnant les valeurs interdites :
\begin{multicols}{2}
\begin{expr}
  \item $2+\dfrac{7}{x+3}$
  \item $4+\dfrac{x-2}{x+1}$
  \item $\dfrac{x}{x-1}-\dfrac{2x}{x+3}$
  \item $\dfrac{4}{x+1}-\dfrac{3}{2x-3}$
  \item $3x-\dfrac{2x+1}{x-2}$
  \item $\dfrac{2x}{x-4}-\dfrac{x+3}{3x-12}$
\end{expr}
\end{multicols}

\exo[Simplifier, multiplier, diviser]
Donner les valeurs interdites, puis simplifier :
\begin{expr}
  \item $\dfrac{4(x+1)^{2}}{8(x+1)}$
  \item $\dfrac{25-x^{2}}{x-5}$
  \item $\dfrac{x+2}{x-2}-\dfrac{x-2}{x+2}$
  \item $\dfrac{x^{2}-1}{x+4}\times\dfrac{x^{2}-16}{x+1}$
  \item $\dfrac{x^{2}-4}{x^{2}+3x}\div\dfrac{x+2}{x+3}$
  \item $\dfrac{2x+1}{x^{2}-9}$
\end{expr}

\partiefiche{Équations quotient}

\exo[Quotient nul]
Donner les valeurs interdites, puis résoudre dans $\R$ :
\begin{multicols}{2}
\begin{enumerate}[label=\alph*)]
  \item $\dfrac{4x-12}{x+2}=0$
  \item $\dfrac{3x+6}{5x^{2}}=0$
  \item $\dfrac{(x-3)(2x+5)}{x+1}=0$
  \item $\dfrac{7-2x}{x^{2}+3}=0$
\end{enumerate}
\end{multicols}

\exo[Se ramener à un quotient nul]
Donner les valeurs interdites, se ramener à un quotient égal à $0$, puis résoudre :
\begin{multicols}{2}
\begin{enumerate}[label=\alph*)]
  \item $\dfrac{3x+2}{x+4}=2$
  \item $\dfrac{-3x}{x+2}=1$
  \item $\dfrac{1}{x}+\dfrac{1}{3}=1$
  \item $\dfrac{5}{x-6}+1=0$
  \item $3-\dfrac{8-x}{x+2}=0$
  \item $\dfrac{x^{2}-16}{x-4}=0$
\end{enumerate}
\end{multicols}

\exo[Égalité de deux quotients]
Donner les valeurs interdites, puis résoudre dans $\R$ :
\begin{enumerate}[label=\alph*)]
  \item $\dfrac{5}{x+2}=\dfrac{3}{x}$
  \item $\dfrac{2x+1}{x}=\dfrac{2x+5}{x+1}$
  \item $\dfrac{3x-2}{x+1}=\dfrac{3x+1}{x-2}$
  \item $\dfrac{x+4}{x-2}-3=\dfrac{8-x}{x-2}$
  \item $\dfrac{x+5}{x-1}=\dfrac{3}{x-1}$
  \item $1-\dfrac{x+1}{x-1}=\dfrac{4}{3-x}$
\end{enumerate}

\partiefiche{Inéquations et tableaux de signes}

\exo[Inéquations du premier degré]
Résoudre dans $\R$ ; donner l'ensemble des solutions sous forme d'intervalle :
\begin{multicols}{2}
\begin{enumerate}[label=\alph*)]
  \item $3x+12>0$
  \item $4x-1<9-x$
  \item $3(x-1)\ge 7-2x$
  \item $\dfrac{2-5x}{3}>0$
  \item $3(x+2)<3x+4$
  \item $\dfrac{x-3}{2}-\dfrac{1-x}{4}>0$
\end{enumerate}
\end{multicols}

\exo[Tableaux de signes]
Dresser le tableau de signes de chaque expression :
\begin{multicols}{2}
\begin{expr}
  \item $(x-5)(x+1)$
  \item $(3-x)(2x+1)$
  \item $x^{2}-16$
  \item $2x(x-3)(4x+1)$
  \item $\dfrac{x+4}{1-x}$
  \item $\dfrac{x(x-2)}{2x+3}$
\end{expr}
\end{multicols}

\exo[Inéquations produit et quotient]
Résoudre dans $\R$, en donnant d'abord les valeurs interdites :
\begin{enumerate}[label=\alph*)]
  \item $\left(9x^{2}-4\right)(x-2)\ge 0$
  \item $\dfrac{5-2x}{x^{2}-4}\le 0$
  \item $\dfrac{x^{2}-9}{3x+1}\ge 0$
  \item $\dfrac{x+3}{2-x}\le 1$
  \item $\dfrac{2-x}{x+4}>3$
  \item $\dfrac{x+1}{x+2}\le\dfrac{x-1}{x+3}$
\end{enumerate}

\partiefiche{Composition de polynômes}

\exo[Premières compositions]
Pour chaque couple $P$, $Q$, écrire $(P\circ Q)(x)$ sous forme développée, réduite et
ordonnée :
\begin{enumerate}[label=\alph*)]
  \item $P(x)=3x+2$ \ et \ $Q(x)=2x+1$
  \item $P(x)=2x+1$ \ et \ $Q(x)=3x+2$
  \item $P(x)=3x-2$ \ et \ $Q(x)=2x^{2}+x-4$
\end{enumerate}
\smallskip
Que remarque-t-on en comparant a) et b) ?

\exo[Composer et calculer]
\begin{enumerate}[label=\arabic*),topsep=0pt]
  \item $P(x)=x^{2}+2x-1$ et $Q(x)=3x-1$ : écrire $(P\circ Q)(x)$.
  \item $P(x)=x^{2}-x$ et $Q(x)=x^{2}+1$ : écrire $(P\circ Q)(x)$.
  \item $P(x)=x^{2}-3$ et $Q(x)=x+4$.
    \begin{enumerate}[label=\alph*)]
      \item Écrire $(P\circ Q)(x)$, puis en déduire $(P\circ Q)(-2)$.
      \item Retrouver ce résultat en calculant $Q(-2)$, puis son image par $P$.
    \end{enumerate}
\end{enumerate}

\exo[Composer dans les deux sens]
\begin{enumerate}[label=\arabic*),topsep=0pt]
  \item $P(x)=x^{2}+1$ et $Q(x)=x^{2}-2x$.
    \begin{enumerate}[label=\alph*)]
      \item Écrire $(P\circ Q)(x)$ et $(Q\circ P)(x)$.
      \item Quel est le degré de $P$, de $Q$, de $P\circ Q$ ?
    \end{enumerate}
  \item $P(x)=2x^{2}-x$ et $Q(x)=\sqrt{2}\,x+1$ : écrire $(P\circ Q)(x)$.
\end{enumerate}

\vfill

\end{multicols}
\end{document}
